\vfill%
\pagebreak
\section{Floating point types}
\label{sec:types:floats}

Numbers with a decimal point are floating-point numbers. Where an integer type
trades away range -- it counts exactly, up to a limit -- a floating-point type
trades away exactness: it reaches enormous magnitudes, and pays for them with
values it can only approximate. How closely it approximates depends on how many
bits it is given.

There are four such types, named on the same pattern as the integers:
\token{f32}, \token{f64}, \token{f80}, and \token{fsize}, whose width is the
widest the machine supports.

\subsection{Floating point literals}

A floating-point literal is two runs of digits with a \token{.} between them and
no spaces anywhere. Underscores may separate digits, as in an integer literal.
Both parts are required: \token{10.0} is a literal, \token{10.} is not.

A literal with no suffix is an \token{f64}. The suffixes are \token{f} for
\token{f32}, \token{d} for \token{f64}, \token{l} for \token{f80} and \token{r}
for \token{fsize} -- the initials of \textit{float}, \textit{double},
\textit{long double} and \textit{real}, the names these widths have carried
since long before Ymir.

\begin{lstlisting}[style=coloredverbatim, label=lst:types:float_literal, caption=Example of floating point literals]
let x = 3.14_15_92f; // f32

let y = 2.71828d; // f64
let z = 2.71828; // f64

let w = 1.0l; // f80

let v = 1_000.0r; // fsize
let u = 1.888_999r; // fsize
\end{lstlisting}

\notebox{
  More advanced forms of floating-point literals are also available,
  such as hexadecimal notation and scientific notation. These representations
  are explained in detail in the advanced chapter on floating-point types (see
  Section~\ref{sec:scalars:floats}).
}

\subsection{Memory representations}

Floating-point numbers follow the IEEE-754 standard, which splits every value
into three fields: a sign, an exponent, and a mantissa -- also called the
significand. How many bits each field gets depends on the type.

\input{chapters/chapter2/figures/float_size}

The value of a floating-point number is computed using the following formula,
where the sign is either $+1$ or $-1$, and the bias is defined as $bias =
2^{e-1} - 1$ with $e$ representing the number of bits used for the exponent. For
instance, in the \token{f32} type, $e = 8$, which gives a bias of $127$.

\begin{equation*}
  value = sign \times mantissa \times 2^{(exponent - bias)}
\end{equation*}

The bias is what lets the exponent go negative without a sign bit of its own: an
exponent of $0$ is stored as $127$, $1$ as $128$, $-1$ as $126$. Positive and
negative exponents end up spread evenly on either side of the stored midpoint,
which keeps the hardware simple. For \token{f32} the exponent runs from about
$-126$ to $127$.

\subsubsection{Example}

Take $87.25$ as an \token{f32}, and work it through field by field.

First, binary. The integer part $87$ is $1010111_2$. The fractional part $0.25$
is $0.01_2$, since $0.25 = 0 \times 2^{-1} + 1 \times 2^{-2} + 0 \times 2^{-3} +
\ldots$ Together they give $1010111.01_2$.

Next, normalisation into the form $1.mantissa \times 2^{exponent}$: shifting the
point six places left gives $1.01011101 \times 2^{6}$. The exponent is therefore
$6$, and since the number is positive the sign bit is $0$.

Then the bias. For \token{f32} it is $127$, so the exponent field holds $6 + 127
= 133$, which on eight bits is $10000101_2$. The mantissa field holds the digits
after the point of $1.01011101$ -- the leading $1$ is implied and not stored --
padded with zeros out to the 23 bits the format allots it.
Table~\ref{tab:types:float_mem_repr_ex} assembles the three.

\begin{table}[h!]
  \centering
  %% The three fields are bit patterns, not quantities: they are set in the
  %% monospaced face like every other bit pattern in the book. As maths they
  %% came out in wide-spaced italics, and the 23-bit mantissa then overran its
  %% column.
  \begin{tabular}{p{.10\textwidth}|p{.22\textwidth}|p{0.42\textwidth}}
    \toprule[0.6pt]
    \textbf{Sign} & \textbf{Exponent} & \textbf{Mantissa} \\
    \toprule[0.6pt]
    \token{0} & \token{10000101} & \token{01011101000000000000000} \\
    \bottomrule[0.2pt]
  \end{tabular}
  \caption{\label{tab:types:float_mem_repr_ex} Memory representation of $87.25$ in a \token{f32}}
\end{table}

\subsection{Stepping}

Between any two decimal numbers there are infinitely many others, and a fixed
number of bits cannot name them all. A floating-point type therefore represents
some values exactly and, for everything in between, silently substitutes the
nearest value it does have. The gaps between the values it has are called
\textit{stepping}.

Listing~\ref{lst:types:stepping_example} shows what that costs. It subtracts
\token{x} from \token{z} to get \token{y}, then adds \token{x} back to \token{y}
to get \token{w}. Arithmetic says \token{w} equals \token{z}. The program says
otherwise: each operation landed on the nearest representable value, and the two
roundings did not cancel. The example uses \token{f32}, and \token{f64} happens
to survive it -- but only this one. Wider types have finer steps, not no steps.

\begin{lstlisting}[style=coloredverbatim, label=lst:types:stepping_example, caption=Example of floating point stepping]
let mut z = 1.989328f; // 'f' suffix to create a f32 literal
let mut x = 0.98769778678f;
let y = z - x;

let w = x + y;

println("(X : ", x, ", Y : ", y, ")");
println("(Z : ", z, ", W : ", w, ")");
println("Z == W : ", z == w);
\end{lstlisting}
\vspace{-10pt}%
\begin{lstlisting}[style=bashVerb]
(X : 0.987698, Y : 1.00163)
(Z : 1.98933, W : 1.98933)
Z == W : false
\end{lstlisting}

\cautionbox{the variable \token{z} and \token{x} are mutable in the
  Listing~\ref{lst:types:stepping_example} to prevent compile time
  optimization and ensure the arithmetic operation are indeed executed at
  runtime.}

\subsection{Special values}

Some bit patterns are reserved. IEEE-754 sets aside particular combinations of
exponent and mantissa to mean infinity, zero, and NaN -- Not a Number, the result
of an operation that has no answer at all.

\input{chapters/chapter2/figures/special_float_numbers}
